\section{Purpose}

\textbf{[STIPULATION]} A190 names the gauge sector as the largest remaining gap.
This document carries out the derivation: from torus topology to the covariant
derivative and the gauge Lagrangian, for all three factors of the SM gauge group
$SU(3)_C\times SU(2)_L\times U(1)_Y$.

\section{Starting Point: Torus and Flux Quantisation}

\textbf{[B]} On $T^4$ there exist non-contractible loops. The flux of a
$U(1)$ connection through such a loop $C$ is quantised: every continuous
deformation of $C$ leaves the flux invariant. This is the basic structure of all
three gauge groups.

\section{$U(1)_Y$: Electromagnetic Charge}

\textbf{[B]} The wave function $\psi$ accumulates a phase upon traversing the
torus loop $C$:
\[
  \psi \;\longrightarrow\; \psi\cdot\exp\!\Bigl(i\,\frac{e}{\hbar}\,\Phi\Bigr),
  \qquad \Phi = \oint_C A_\mu\,dx^\mu = n\,\frac{h}{e},\quad n\in\mathbb{Z}.
\]
The monodromy $\exp(i\cdot 2\pi n)=1$ is satisfied for all $n\in\mathbb{Z}$.
This is the $U(1)$ gauge structure: the phase $e^{i\alpha(x)}$ with
$\alpha(x) = (e/\hbar)\int A_\mu dx^\mu$.

\textbf{[B]} The covariant derivative and the field-strength tensor follow:
\[
  D_\mu = \partial_\mu + ieA_\mu,
  \qquad
  F_{\mu\nu} = \partial_\mu A_\nu - \partial_\nu A_\mu,
  \qquad
  \mathcal{L}_\text{EM} = -\tfrac14\,F_{\mu\nu}F^{\mu\nu}.
\]
The gauge Lagrangian follows from minimal coupling: $F_{\mu\nu}$ is the curvature
of the connection $A_\mu$, $\mathcal{L}_\text{EM}$ is its norm. No free step.

\section{$SU(3)_C$: Colour Charge from Linking Number}

\textbf{[B]} Three flux tubes on $T^4/\mathbb{Z}_3$ carry linking numbers
$\ell k(i,j)\in\mathbb{Z}$ for $i,j\in\{1,2,3\}$ (tubes $i$ and $j$).
The independent non-trivial states are:

\begin{center}
\renewcommand{\arraystretch}{1.3}
{\small\begin{tabular}{lll}
\toprule
Linking state & Generator & Physics \\
\midrule
$\ell k(1,2)=\pm1$, rest~0 & $\lambda_1,\lambda_2$ & Gluon $r\leftrightarrow g$ \\
$\ell k(1,3)=\pm1$, rest~0 & $\lambda_4,\lambda_5$ & Gluon $r\leftrightarrow b$ \\
$\ell k(2,3)=\pm1$, rest~0 & $\lambda_6,\lambda_7$ & Gluon $g\leftrightarrow b$ \\
$\ell k(1,1)-\ell k(2,2)$  & $\lambda_3$           & Gluon neutral (1) \\
$\ell k(1,1)+\ell k(2,2)-2\,\ell k(3,3)$ & $\lambda_8$ & Gluon neutral (2) \\
\bottomrule
\end{tabular}}
\renewcommand{\arraystretch}{1.0}
\end{center}

\textbf{[B]} Count: $\binom{3}{2}=3$ pairs, each with real and imaginary part,
gives $6$ off-diagonal; plus $2$ traceless diagonal generators (rank-$1$ reduction
of the $3\times3$ tracelessness condition) gives $6+2=8$ independent generators
-- exactly $\dim\bigl(su(3)\bigr)$. The gauge group $SU(3)_C$ is topologically
forced, not postulated.

\textbf{[B]} Covariant derivative and Lagrangian analogous to $U(1)$:
\[
  D_\mu = \partial_\mu + ig_s\,G_\mu^a\,\lambda_a/2,
  \qquad
  G_{\mu\nu}^a = \partial_\mu G_\nu^a - \partial_\nu G_\mu^a
    + g_s\,f^{abc}\,G_\mu^b\,G_\nu^c,
\]
\[
  \mathcal{L}_\text{strong} = -\tfrac14\,G_{\mu\nu}^a\,G^{a,\mu\nu}.
\]
The non-commutativity $f^{abc}\neq0$ follows from the non-commutativity of the
linking-number operations for three tubes.

\section{$SU(2)_L$: Weak Isospin Group from Chiral Projection}

\textbf{[B]} The chiral projection $P_+$ (A095) selects the $m\ge0$ sector
of the torus. The spinor-doublet character of the weak interaction follows
from the spin-$\tfrac12$ structure (winding $(1,1)$, A263): the two
components of the isospin doublet $(u_L, d_L)$ and $(\nu_L, e_L)$ respectively
are the two independent $m=0$ and $m=+1$ sectors after projection.

\textbf{[B]} The three generators $\tau^+, \tau^-, \tau^3$ (Pauli matrices)
of $SU(2)_L$ correspond to the three independent transitions between the
two projected states. Covariant derivative:
\[
  D_\mu = \partial_\mu + ig_w\,W_\mu^a\,\sigma_a/2,
  \qquad
  W_{\mu\nu}^a = \partial_\mu W_\nu^a - \partial_\nu W_\mu^a
    + g_w\,\varepsilon^{abc}\,W_\mu^b\,W_\nu^c,
\]
\[
  \mathcal{L}_\text{weak} = -\tfrac14\,W_{\mu\nu}^a\,W^{a,\mu\nu}.
\]

\section{Overall Picture: Gauge Lagrangian from Torus Topology}

\textbf{[B]} The complete Standard Model gauge Lagrangian:
\[
  \mathcal{L}_\text{gauge}
  = -\tfrac14\,G_{\mu\nu}^a G^{a,\mu\nu}
  - \tfrac14\,W_{\mu\nu}^a W^{a,\mu\nu}
  - \tfrac14\,F_{\mu\nu}F^{\mu\nu}
\]
follows from three topological structures of $T^4/\mathbb{Z}_3$:
\begin{itemize}[leftmargin=2em,itemsep=2pt,topsep=2pt]
  \item $U(1)_Y$: flux quantisation through non-contractible loop
  \item $SU(3)_C$: linking number of three $\mathbb{Z}_3$ flux tubes
    ($6+2=8$ generators)
  \item $SU(2)_L$: chiral projection of the spin-$\tfrac12$ sector (A095)
\end{itemize}

\begin{center}
\renewcommand{\arraystretch}{1.3}
{\small\begin{tabular}{lccc}
\toprule
Group & Topological origin & Generators & Gauge bosons \\
\midrule
$U(1)_Y$    & Flux quantisation  & 1 & Photon ($\gamma$) \\
$SU(2)_L$   & Chiral projection  & 3 & $W^+,W^-,Z^0$ \\
$SU(3)_C$   & Linking number     & 8 & 8 gluons \\
\midrule
Total & & 12 & 12 SM gauge bosons \\
\bottomrule
\end{tabular}}
\renewcommand{\arraystretch}{1.0}
\end{center}

\section{Symbol Glossary}

Symbols used throughout the A-Series are explained in A010.
Specific to this document:

\begin{center}
\renewcommand{\arraystretch}{1.3}
{\small\begin{tabular}{lp{9cm}}
\toprule
Symbol & Meaning \\
\midrule
$\ell k(i,j)$ & Linking number of tubes $i$ and $j$ on $T^4/\mathbb{Z}_3$ \\
$\lambda_a$ & Gell-Mann matrices, generators of $SU(3)$ \\
$\sigma_a$ & Pauli matrices, generators of $SU(2)$ \\
$G_\mu^a, W_\mu^a, A_\mu$ & Gauge fields of $SU(3)$, $SU(2)$, $U(1)$ \\
$f^{abc}$ & Structure constants of $SU(3)$ \\
$\varepsilon^{abc}$ & Levi-Civita symbol, structure constants of $SU(2)$ \\
$g_s, g_w, e$ & Coupling constants (not derived from $\xi$) \\
\bottomrule
\end{tabular}}
\renewcommand{\arraystretch}{1.0}
\end{center}

\section{Verification Script}

\begin{itemize}[leftmargin=2em,itemsep=2pt,topsep=2pt]
  \item Count of linking states; Gell-Mann matrix identification;
    tracelessness condition; $6+2=8$:
    \path{python/A_Serie_Skripte/a192_eichsektor.py}
\end{itemize}

\section{Open Questions}

\textbf{The coupling constants $g_s$, $g_w$, $e$ are not derived from $\xi$.}
The group structure follows from the topology; the strength of the couplings
does not. This is the remaining gap.

\textbf{The running of the couplings} ($g_s(\mu)$, $g_w(\mu)$ as functions of
the energy scale $\mu$) is not derived from FFGFT.

\textbf{The Higgs mechanism} -- the spontaneous symmetry breaking
$SU(2)_L\times U(1)_Y\to U(1)_\text{EM}$ -- is not derived from the
torus geometry. The Higgs-EFT consistency check (A190, $2{.}3\,\%$) is a
hint, not a proof.

\section{Supersedes}

This document systematises the gauge-sector content of the old corpus
(charge quantisation and colour charge from torus topology, Lagrangian structure)
self-contained without old-corpus citations. It supplements A190 (SM as
projection), A095 (chiral projection) and A075 (field theory).

\section{Use of AI Tools}

This document was created with the assistance of AI tools.
Responsibility for content and errors rests with the author.
Full wording of the rules in A010.
